\chapter{Physical Mathematics}

OmniLaTeX provides a comprehensive physical typesetting system built on
\texttt{siunitx}~v3 for units, \texttt{chemmacros} for chemistry, and
\texttt{xfrac}/\texttt{nicefrac} for fractions.  All configuration is
handled by \texttt{omnilatex-math.sty}, so you can use these commands
directly without additional setup.

% =========================================================================
\section{siunitx Setup}
% =========================================================================

\texttt{siunitx} is preconfigured in \texttt{omnilatex-math.sty} with
locale-aware settings.  The base configuration applies to all languages,
and per-language overrides are activated automatically based on the
document language:

\begin{verbatim}
\sisetup{%
    mode = match,
    per-mode = symbol,
    range-units = single,
    uncertainty-mode = separate,
    propagate-math-font = true,
    ...
}
\end{verbatim}

When the document language is German, the locale is set to \texttt{DE}
(uses a comma as decimal separator).  For English documents the locale
defaults to \texttt{US} (period as decimal separator).  This switching
happens via \cs{gappto} and requires no manual intervention.

% =========================================================================
\section{Core Commands}
% =========================================================================

The three primary \texttt{siunitx} commands cover most use cases:

\begin{table}[htbp]
  \centering
  \caption{Core \texttt{siunitx} commands.}
  \label{tab:siunitx-core}
  \begin{tabular}{@{}lp{0.55\textwidth}@{}}
    \toprule
    Command & Description \\
    \midrule
    \texttt{\textbackslash num\{value\}} & Formats a plain number
      respecting locale (e.g.\ \texttt{1,000.5} in US,
      \texttt{1.000,5} in DE). \\[4pt]
    \texttt{\textbackslash unit\{...\}} & Typesets a unit with correct
      spacing and symbols, e.g.\
      \texttt{\textbackslash unit\{\textbackslash metre\textbackslash per\textbackslash second\}}
      $\to$ \unit{\metre\per\second}. \\[4pt]
    \texttt{\textbackslash qty\{value\}\{unit\}} & Combines number and unit
      in one call, e.g.\
      \texttt{\textbackslash qty\{5.2\}\{\textbackslash kilo\textbackslash gram\}}
      $\to$ \qty{5.2}{\kilo\gram}. \\
    \bottomrule
  \end{tabular}
\end{table}

Additional useful commands include \texttt{\textbackslash qtylist} for
comma-separated quantities (e.g.\
\texttt{\textbackslash qtylist\{1;2;3\}\{\textbackslash metre\}}) and
\texttt{\textbackslash numlist} for plain number lists.

% =========================================================================
\section{Quantity Ranges}
% =========================================================================

Ranges are typeset with \texttt{\textbackslash qtyrange}, which by default
places the unit at the end once (\texttt{range-units=single}):

\begin{verbatim}
\qtyrange{20}{30}{\celsius}
% → 20 °C to 30 °C  (en-dash separator)
\end{verbatim}

The range phrase can be customised:

\begin{verbatim}
\qtyrange[range-phrase={/}]{20}{30}{\celsius}
% → 20/30 °C
\end{verbatim}

OmniLaTeX's \texttt{\textbackslash temperaturepair} macro (see
Section~\ref{sec:temp-pair}) is a specialised wrapper that uses the
\texttt{/} separator for temperature ranges.

% =========================================================================
\section{Custom SI Units}
% =========================================================================

OmniLaTeX declares several additional units beyond the \texttt{siunitx}
defaults.  These are defined in \texttt{omnilatex-math.sty} and are
available immediately:

\begin{table}[htbp]
  \centering
  \caption{Custom SI units defined by OmniLaTeX.}
  \label{tab:custom-units}
  \begin{tabular}{@{}llp{0.45\textwidth}@{}}
    \toprule
    Command & Unit & Description \\
    \midrule
    \texttt{\textbackslash volpercent}       & vol\%  & Volume percent \\
    \texttt{\textbackslash watthour}         & Wh     & Watt-hour \\
    \texttt{\textbackslash annum}            & a      & Annum (year) \\
    \texttt{\textbackslash atmosphere}       & atm    & Standard atmosphere \\
    \texttt{\textbackslash partspermillion}  & ppm    & Parts per million \\
    \texttt{\textbackslash bar}              & bar    & Bar (pressure) \\
    \texttt{\textbackslash relhumidity}      & \%rh   & Relative humidity (language-dependent) \\
    \bottomrule
  \end{tabular}
\end{table}

The relative humidity unit adapts to the document language: it renders as
\texttt{\%\,r.F.} in German and \texttt{\%\,RH} in English.

The Euro sign is also available as a unit via the \texttt{eurosym}
package---use \texttt{\textbackslash euro} directly inside
\texttt{\textbackslash unit\{\}} or \texttt{\textbackslash qty\{\}\{\}}.

% =========================================================================
\section{Custom SI Qualifiers}
% =========================================================================

Unit qualifiers appear as subscripts and are used to disambiguate
quantities that share the same unit.  OmniLaTeX integrates qualifiers
with the glossary system so that subscript labels remain consistent
throughout the document:

\begin{table}[htbp]
  \centering
  \caption{Custom SI qualifiers defined by OmniLaTeX.}
  \label{tab:si-qualifiers}
  \begin{tabular}{@{}llp{0.45\textwidth}@{}}
    \toprule
    Command & Qualifier & Description \\
    \midrule
    \texttt{\textbackslash dryair}  & dry air  & Dry air component \\
    \texttt{\textbackslash water}   & water    & Water component \\
    \texttt{\textbackslash thermal} & thermal  & Thermal property \\
    \bottomrule
  \end{tabular}
\end{table}

Usage example:

\begin{verbatim}
\unit{\kilogram\per\cubic\metre\dryair}
% → kg/m³_dry air
\end{verbatim}

The qualifier text is sourced from the glossary entries
\texttt{sub.dry\textasciitilde air}, \texttt{sub.water}, and
\texttt{sub.thermal}, ensuring consistent subscript formatting.

% =========================================================================
\section{Chemistry}
% =========================================================================

Chemical typesetting is provided by the \texttt{chemmacros} package,
configured to use the Libertinus font family for formulae.

\textbf{Compounds} are typeset with \texttt{\textbackslash chcpd}:

\begin{verbatim}
\chcpd{H2O}    % → H₂O (compound style)
\ch{H2O}       % → H₂O (formula style)
\ch{CO2 + H2O -> H2CO3}
\end{verbatim}

\textbf{Numbered reactions} use the \texttt{reaction} environment:

\begin{verbatim}
\begin{reaction}
  2 H2 + O2 -> 2 H2O
\end{reaction}
% → R [1]:  2 H₂ + O₂ → 2 H₂O
\end{verbatim}

Reactions are tagged with a bold \textbf{R} prefix and square brackets,
sharing the equation counter.  In chaptered documents the counter is
prefixed with the chapter number.

\textbf{Chemical amounts} combine quantity, unit, and compound:

\begin{verbatim}
\chemamount{1.2}{\partspermillion}{CO2}
% → 1.2 ppm CO₂
\end{verbatim}

This is a convenience macro defined in \texttt{omnilatex-math.sty} that
ensures correct spacing between the numerical value, unit, and compound
name.

\textbf{Multiple reactions} can be grouped:

\begin{verbatim}
\begin{reactions}
  A -> B & C -> D \\
  B + D -> E
\end{reactions}
\end{verbatim}

A gathered variant without alignment is also available:
\texttt{reactionsgather}.

% =========================================================================
\section{Fractions}
% =========================================================================

Two packages provide alternative fraction styles:

\textbf{xfrac} --- stacked (slanted) fractions for inline use:

\begin{verbatim}
\sfrac{1}{2}  % → ¹⁄₂  (diagonal fraction)
\end{verbatim}

\textbf{nicefrac} --- inline fractions with a slash:

\begin{verbatim}
\nicefrac{1}{3}  % → 1/3  (compact inline fraction)
\end{verbatim}

Both packages are loaded by \texttt{omnilatex-math.sty}.  Use
\texttt{\textbackslash sfrac} when you want a true diagonal fraction
 glyph, and \texttt{\textbackslash nicefrac} when a simple slash is
 preferable for very compact inline settings.
